\chapter{消费组、重平衡与提交围栏}
消费组是一群协作读取同一 topic 的 consumer：成员是参与者标识，assignment 是当前每个成员拥有的分区列表，同组一个分区同时只分给一个成员。负责保存组状态并答复加入/查询的 coordinator 在成员变化时重算整个 assignment，这次变化叫 rebalance。generation 是每次拓扑变化递增的组版本，不是时钟或消费 offset；稍后的 GroupToken 会把组、成员和 generation 一起带入请求，使服务端能拒绝旧成员。心跳用于表示成员仍活跃，过期成员会被移除。本课程是同一 topic、即时重分配、单进程可信 coordinator 的教学模型，不是完整 Kafka 组协议或高可用共识。

\begin{figure}[ht]
\centering
\begin{tikzpicture}[font=\small,node distance=10mm]
\node[draw,rounded corners,minimum width=19mm] (a1) {A joins};
\node[draw,rounded corners,minimum width=26mm,right=of a1] (g1) {generation 1};
\node[draw,rounded corners,minimum width=19mm,right=of g1] (b) {B joins};
\node[draw,rounded corners,minimum width=26mm,right=of b] (g2) {generation 2};
\node[draw,rounded corners,minimum width=23mm,below=of g2] (fence) {A gen 1 fenced};
\draw[-{Stealth},thick] (a1)--(g1); \draw[-{Stealth},thick] (g1)--(b);
\draw[-{Stealth},thick] (b)--(g2); \draw[-{Stealth},thick] (g2)--(fence);
\end{tikzpicture}
\caption{成员拓扑变化递增 generation；旧 token 不能覆盖当前提交。}
\end{figure}

\lessonsection{17}{分区分配}
\goals{实现确定、平衡且每个分区只有一个 owner 的 Round Robin 分配。}{理解去重、稳定排序、不可变结果和商余数分配。}
\background{分区可以在 consumer 间并行处理，但若同一分区被两名成员同时处理，进度会冲突。Round Robin 在排序后的成员间循环派发排序后的分区；结果可重现，每个成员分到的分区数最多相差一。这里的 assignor 是纯计算函数，尚不负责判断成员是否真的加入了 group。}
\subsubsection*{开始前先读}
先复习 \stepref{3}{TopicPartition 与分区排序}。本步的分配器会被 \stepref{18}{成员 coordinator} 调用；它只根据传入的 partition/member 列表分配，不查询目录、不改日志。
\providededit{\method{TopicPartition} 值类型、group 返回值类型已提供；\method{RoundRobinAssignor} 构造器无状态，不提供分配算法答案。}{\pathcode{exercises/src/main/java/io/simplekafka/group/RoundRobinAssignor.java}：\method{Map<String,List<TopicPartition>> assign(List<TopicPartition> partitions, List<String> members)}。只实现 \method{assign}。}
\subsubsection*{参数、结果与状态}
\begin{itemize}
\item partitions 是待分配的 \pathcode{TopicPartition} 列表，members 是成员标识列表；两列表及元素均不可为 null。相同 partition/member 只保留一份。
\item partition 按 topic 后 partition number 升序，member 按 Java \pathcode{String} 自然顺序排序；不是输入顺序，也不是数字字符串的数值顺序。
\item 返回 map 的每个不同 member 都有列表，包括空列表；map 和内部列表均不可变。assignor 不改变输入列表，也没有跨调用状态。
\item null 输入/元素，或 member 标识 UTF-8 编码长度不在 1--255 bytes，报 \method{INVALID_REQUEST}。无 member 时返回空 map；有 member 但无 partition 时保留所有成员并给空列表。
\end{itemize}
\lessoncontract{至少有一个不同 member 时，每个不同 partition 恰好出现一次且 owner 属于 members；所有不同 member 都出现在结果中。成员分区数量最大值与最小值之差不超过 1。无 member 时遵循上项，返回空 map，不产生 owner。}
\expected{初始 partitions 明确为 \pathcode{[orders-0, orders-1, orders-2, orders-3, orders-4]}，members 输入顺序为 \pathcode{[b,a,a]}。排序去重后 members=[a,b]，分区 0、2、4 给 a，1、3 给 b；改乱输入顺序结果不变。若有 5 个分区和按 Java 字符串顺序排列的 \pathcode{[a,b,c,d,e,f,g]}，前五人各得一个、f/g 保留空列表。\\
\method{Step17Test.sortsAndDeduplicatesPartitionsAndMembersBeforeRoundRobinAssignment} 检查去重与确定顺序；\method{assignsFivePartitionsToTheFirstFiveOfSevenSortedMembersAndKeepsEmptyMembers} 检查多成员空列表；\method{sortsCrossTopicPartitionsAndMemberIdsLexicallyAfterDeduplication} 检查跨 topic 与 \pathcode{"10"} 排在 \pathcode{"2"} 前的词典顺序。}
\lessonhints{把“唯一 owner”拆成完整覆盖和没有重复两项分别检查。}{先排好成员与分区，再按索引取模分配；输入 List 的顺序不是本课程规则。}{分别手算 5 分区/2 成员与 5 分区/7 成员，并把重复输入加入测试，确认它不会让一个分区出现两次。}
\lessonquestions{为什么稳定排序是可观察契约，而不只是输出美观？}{成员多于分区时，为什么仍要返回空列表成员？}
\paragraph{自检解释。}稳定排序让相同集合每次产生相同 owner，便于协调与验证；空列表明确表达成员仍在组中但当前无分区。
\pitfalls{只返回有分区成员会丢失空 assignment；依赖偶然 map 遍历顺序会破坏可重现性。真实 Kafka 有多种 assignor 和协作式重平衡，本课程只做单 topic 的简单 Round Robin，不声称它是 Kafka 默认算法。}
\lessoncommand{17}

\lessonsection{18}{成员与重平衡}
\goals{管理 group 成员、generation 与全组分配，在成员拓扑变化时正确重平衡。}{理解 join、leave、generation 与 committed offset 各自代表什么。}
\background{成员是以 member id 标识的 consumer 实例；一个 group 在本课程中订阅一个 topic。coordinator 保存成员集合及它们共享的 generation，并用第 17 步 assignor 计算完整分配。新成员加入或成员离开改变拓扑，generation 加一并立即重算所有成员的分区。重复同一 member/topic 的 join 只刷新心跳，不改变拓扑；generation 不是时间戳，也不保存消费进度。}
\subsubsection*{开始前先读}
先掌握 \stepref{10}{topic metadata 与分区目录}，再看 \stepref{17}{确定分配}。需要知道一个 topic 对应哪些 \pathcode{TopicPartition}，并区分 coordinator 的内存成员状态与第 15 步的持久消费位点。
\providededit{\method{PartitionCatalog}、组锁、成员状态容器、\method{TimeSource} 与第 17 步 assignor 已提供；消息请求值类型已编码。}{\pathcode{exercises/src/main/java/io/simplekafka/group/GroupCoordinator.java}：\method{GroupAssignment join(String group, String member, String topic)}、\method{GroupAssignment leave(String group, String member)}、\method{GroupAssignment assignment(String group, String member)}；并在 \pathcode{exercises/src/main/java/io/simplekafka/broker/BrokerHandler.java} 的 \method{handle(Messages.Request)} 增补 \pathcode{JOIN_GROUP/LEAVE_GROUP/GROUP_ASSIGNMENT}。}
\subsubsection*{参数、结果与状态}
\begin{itemize}
\item group/member 必须是非空且不超过 255 UTF-8 bytes 的标识；topic 必须是有效课程 topic。首次 join 只在 topic 已存在后创建组，初始 generation=0，成功加入后为 1。
\item 新成员加入或已知成员成功离开，generation 恰加一，返回包含全体当前成员的 \method{GroupAssignment}。同 member、同 topic 重复 join 仅更新 lastHeartbeat 并返回现有 assignment；不得重复分配或增加 generation。
\item 已存在 group 的 topic 冲突优先报 \method{INVALID_REQUEST}；未知 topic 报 \method{UNKNOWN_TOPIC_OR_PARTITION}，且不创建 group/消耗 generation。标识无效报 \method{INVALID_REQUEST}。未知 group/member 的 leave 或 assignment 报 \method{UNKNOWN_MEMBER}；coordinator 关闭后操作报 \pathcode{IllegalStateException}。
\item 组变空时保留 coordinator 内的 generation 历史和已提交 offsets；offset store 与成员生命周期分离。成员集合本身是进程内状态，不因 offsets 持久而变成持久组协议。
\end{itemize}
\lessoncontract{每次真实拓扑变化只增加一次 generation，并为全部剩余成员重新计算分配；重复 join 不改变 generation。assignment 查询只对当前已知成员成功，且返回同一份全组映射。}
\expected{初始 \pathcode{workers} 组不存在，\pathcode{orders} 有分区 0、1、2。a join 得 generation 1、拥有全部三分区；a 同 topic 再 join 仍为 generation 1。b join 后 generation 2，a 得 0、2，b 得 1；c join 后 generation 3，三人各得一个。c leave 后 generation 4，重算为 a 得 0、2、b 得 1。若随后组变空，generation 仍保留，原有 group offsets 仍可 fetch。\\
\method{Step18Test.membershipChangesRebalanceEverySurvivorAndRetainEmptyGroupGenerationOverTcp} 通过真实 TCP 检查加入、重复加入、全组重算、topic/unknown-member 错误、空组进度保留及再次加入。}
\lessonhints{先判断请求是否真的改变成员集合；重复 join、被拒绝的 join 与新成员加入不是同一状态变化。}{每次拓扑变化后检查所有成员都看到相同 generation 和完整 assignment，不能只更新新成员。}{从零列出 a 加入、a 重复、b 加入、c 加入、c 离开的 generation 和每个 owner。}
\lessonquestions{为什么 b 加入时 a 也必须获得新 generation 的 assignment？}{最后一个成员离开后，为什么 group offset 不该删除？}
\paragraph{自检解释。}分区可能已从 a 转交给 b，若 a 保留旧 assignment 就无法对齐服务端当前所有权；consumer 进度属于 group/topic-partition，不属于仍在线的 member，所以空组也可保留进度。
\pitfalls{给每个成员独立递增版本会使 token 不一致；组空不等于消费进度失效。真实 Kafka 有加入、同步 assignment 与协作撤销等多阶段协议；本课程即时重分配，不模拟完整 classic barrier 或 KIP-848。}
\lessoncommand{18}

\lessonsection{19}{心跳与过期}
\goals{用可控时钟识别失联成员，并把同组到期成员合并为一次拓扑变化。}{理解 lastHeartbeat、session timeout 和包含等号的边界。}
\background{成员进程可能崩溃而来不及 leave；coordinator 不能只凭沉默立刻移除它。成功 heartbeat 记录注入式时钟当前毫秒值，表示该成员仍活跃。若经过的时间达到 session timeout，成员过期。确定性时钟让边界测试不需依赖 sleep；提供的周期 worker 是运行时触发机制，和测试直接调用过期 hook 是两回事。}
\subsubsection*{开始前先读}
需要 \stepref{18}{成员、generation 和 assignment}。要能解释同一 group 的成员共享 generation，并知道被移除成员的旧 token 不再是有效心跳。
\providededit{\method{TimeSource}、组锁、正 session timeout 和周期调度生命周期已提供。}{\pathcode{exercises/src/main/java/io/simplekafka/group/GroupCoordinator.java}：\method{void heartbeat(GroupToken token)}、\method{Set<String> expire()}；并在 \pathcode{exercises/src/main/java/io/simplekafka/broker/BrokerHandler.java} 的 \method{handle(Messages.Request)} 增补 \pathcode{HEARTBEAT}。}
\subsubsection*{参数、结果与状态}
\begin{itemize}
\item coordinator 构造时 sessionTimeoutMillis 必须大于 0，否则 \method{INVALID_REQUEST}。heartbeat token 为 null 报 \method{INVALID_REQUEST}；先判断 member 是否存在，不存在报 \method{UNKNOWN_MEMBER}；之后才检查 generation，旧值或 future 值报 \method{ILLEGAL_GENERATION}。失败不刷新时间；有效 token 才把该 member 的 lastHeartbeat 更新到当前 clock。coordinator 关闭后操作报 \pathcode{IllegalStateException}。
\item \method{expire()} 不接收参数，使用注入的当前时钟，返回已移除成员的排序集合，元素格式为 \pathcode{group/member}。当 \pathcode{now-lastHeartbeat >= sessionTimeoutMillis} 时移除；未达到阈值不变，时间戳位于未来的成员也不会提前过期。
\item 一次 expire 同一 group 中可同时移除多个到期成员，但该 group generation 只增加一次、只重算一次；其余未到期成员保留。返回空集合表示本轮没有移除成员；重复调用无新变化。过期不会删 committed offsets。
\item broker 的周期 worker 会自动触发检查；确定性测试应移动 \method{TimeSource} 再调用提供的 \pathcode{expireGroups()} 或 coordinator \method{expire()}，而不是用真实 sleep 猜时序。
\end{itemize}
\lessoncontract{过期边界采用 \pathcode{>=}；无效/旧 heartbeat 不续期。过期结果按 group 再 member 升序排列。每个因到期而变化的 group 每轮 generation 仅增加一次；组空后仍保留 generation。}
\expected{设 clock=1000、timeout=100，a 与 b 均在 1000 加入 \pathcode{workers}，此时 generation=2，a 的 generation-1 heartbeat 已过期。clock=1099 时，a 用旧 token 或 future generation 99 心跳均报 \method{ILLEGAL_GENERATION}，不存在的成员报 \method{UNKNOWN_MEMBER}；这些失败不刷新。b 用 generation 2 成功 heartbeat，更新为 1099。1099 调 expire 返回空；clock=1100 再调，只移除 \pathcode{workers/a}（恰好经过 100 ms），generation 变 3，b 获得全部三分区。\\
\method{Step19Test.rejectsStaleAndFutureHeartbeatsAndExpiresOnlyTheUnrefreshedMemberAtTheBoundary} 检查 1099/1100 与错误 heartbeat；\method{expirationOrdersAllIdentifiersAndAdvancesEachChangedGroupOnce} 检查排序及同组一次递增；\method{idempotentJoinRefreshesHeartbeatButRejectedJoinDoesNot} 检查续期；\method{replicatedClusterExpiresIdleMembersWithoutIncomingRequests} 与 \method{expireGroupsProvidesDeterministicClusterExpiryAndChecksLifecycle} 分别检查周期 worker 和确定性 hook。}
\lessonhints{在同一状态观察中使用当前时钟、判断 elapsed 时间并移除成员；失败 heartbeat 不应写 lastHeartbeat。}{返回值要能指出具体的 group/member，而不仅是过期数量。}{timeout=100、lastHeartbeat=1000 时，比较 1099 与 1100；再在 1099 刷新 B，确认本轮只移除 A。}
\lessonquestions{为什么必须在契约中明确 timeout 的等号边界？}{同组 A、C 在一次 expire 中都到期，为什么 generation 不能加两次？}
\paragraph{自检解释。}若不规定等号，恰好到期时就会产生不同的成员所有权结果；一次 expire 是一次观察到的拓扑变化，所以同组一起移除只触发一次全组重分配。
\pitfalls{用系统时间和 sleep 测精确阈值会受调度影响；边界用注入 clock，worker 另做集成检查。先写 lastHeartbeat 再验证 token 会让过期成员非法续命。真实 Kafka 心跳还受配置、网络和协调协议影响。}
\lessoncommand{19}

\lessonsection{20}{消费组围栏}
\goals{将 assignment 与 fetch/offset commit 结合，用 generation token 拒绝僵尸成员的过期操作。}{理解 token 三元组、撤销分区、刷新 assignment 与外部副作用边界。}
\background{rebalance 后，旧 consumer 线程可能仍在运行。GroupToken=(group,member,generation) 是请求携带的组成员版本；broker 必须确认成员仍存在、generation 当前、该成员拥有目标分区，才允许带 token 的 FETCH 或 COMMIT。token 围栏能拒绝旧请求，但不能撤回已经送到旧客户端的记录，也不能撤销已发生的业务副作用。手动模式不带 token，仍由此前的普通 consumer 路径使用。}
\subsubsection*{开始前先读}
需要 \stepref{14}{fetch 的 position}、\stepref{15}{持久 offset key}、\stepref{17}{分配算法}、\stepref{18}{generation/assignment}及 \stepref{19}{heartbeat/expiry}。本步把它们组合起来；不要把成员 generation 当成消费 offset。
\providededit{group 请求/响应类型、\method{RpcClient}、\method{SimpleConsumer} 的 pull/offset 路径及成员生命周期骨架已提供。}{\pathcode{exercises/src/main/java/io/simplekafka/group/GroupCoordinator.java}：\method{void validate(GroupToken token, TopicPartition tp)}；\pathcode{exercises/src/main/java/io/simplekafka/client/GroupConsumer.java}：\method{GroupConsumer(RpcClient client, String group, String member)}（接管 client 并在 close 时关闭）、\method{GroupConsumer(MetadataRouter router, String group, String member)}（借用 router，不关闭它）；两个构造器的 group/member 均须为 1--255 UTF-8 bytes。只实现 \method{subscribe(String topic)}、\method{Map<TopicPartition,List<LogRecord>> poll(int maxRecords, int maxBytes)}、\method{commitSync()}；并在 \pathcode{exercises/src/main/java/io/simplekafka/broker/BrokerHandler.java} 增补带 token 的 \pathcode{FETCH/COMMIT} 围栏。}
\subsubsection*{参数、结果与状态}
\begin{itemize}
\item \method{GroupCoordinator.validate(token,tp)} 要求当前存在的 member、当前 generation 和该 member 对目标分区的所有权。未知 member 报 \method{UNKNOWN_MEMBER}，generation 不符（包括 future token）报 \method{ILLEGAL_GENERATION}，未分配给此 member 报 \method{NOT_ASSIGNED}；null tp 报 \method{INVALID_REQUEST}，拒绝不会写 OffsetStore。
\item \method{subscribe(topic)} 显式 join 并拿到 assignment/token；新分区从 committed offset 或 0 恢复，仍保留的分区保留原 position，撤销分区丢弃。未订阅就 poll/commit 报 \method{INVALID_REQUEST}；未知 topic 等错误向调用者传播。consumer 关闭后操作报 \pathcode{IllegalStateException}。
\item \method{poll(maxRecords,maxBytes)} 要求两个总预算均为正，否则报 \method{INVALID_REQUEST}，并在 fetch 请求前 heartbeat；返回的 map 不可变。若 heartbeat 请求返回 \method{ILLEGAL_GENERATION}，只查询当前 assignment、更新本地 token/position，再用新 token fetch；撤销位置删除，新分区从 committed/0 起。遇 \method{UNKNOWN_MEMBER} 不自动 join，必须由调用者显式 subscribe；其他错误传播。成功 fetch 推进相应的内存 position，不直接提交。
\item \method{commitSync()} 用当前 token 对每个已分配分区分别提交 position；它不 heartbeat，也不自行刷新 assignment。旧 token/非 owner 被拒时 offset store 不变；逐分区请求不是原子事务，较早提交可能在后续请求失败前已成功。broker 的 COMMIT 分支要求非负 nextOffset，且 OffsetKey 的 group 与 token、分区与请求一致；否则 \method{INVALID_REQUEST} 且不追加。服务端在 coordinator 锁内完成 validate 与 store append，锁顺序 coordinator 再 store；它们不是分区日志事务。普通手动请求的 null token 仍合法。
\item 协议成功时 token FETCH 返回 \pathcode{FetchBody}，token COMMIT 返回 \pathcode{EmptyBody}；assignment 查询返回 \pathcode{GroupBody}。已发送记录和客户端副作用不因后续拒绝而回滚。coordinator 关闭后操作报 \pathcode{IllegalStateException}。
\end{itemize}
\lessoncontract{FETCH 与 COMMIT 都围栏 member、generation 和分区 owner。COMMIT 的 OffsetKey 必须与 token 的 group 及请求分区匹配；任何拒绝不得追加该次 offset。generation 变化后，旧成员不能提交已撤销分区的缓存 position。}
\expected{初始 \pathcode{orders-0} 只有 offset 0；\pathcode{orders-1} 有 offset 0--8，且 \pathcode{workers/orders-1} 已提交 7。a subscribe 得 generation 1、分到两分区；poll(1,4096) 先读 p0/offset0，position 为 p0=1、p1=7。b subscribe 后 generation 2，a 只持有 p0、b 持有 p1。a 用旧 token commitSync 报 \method{ILLEGAL_GENERATION}，p0 仍未提交、p1 仍为 7。此时 p0 追加 offset1；a poll 检测旧 generation，获取新 assignment，保留 p0=1、丢弃已撤销 p1，再读 p0/offset1；commit 后 p0=2，p1 仍为 7。b 首次 poll 从已提交的 p1=7 开始。若 broker fixture 重启，成员内存状态消失、日志和 offsets 保留（p0=2、p1=9）；创建新的 \method{GroupConsumer} a 并显式 subscribe，generation 从 1 开始，新 position 从已提交的 2、9 恢复。\\
\method{Step20Test.coordinatorFencesCommitsByGenerationOwnershipAndMatchingOffsetKey} 检查 coordinator 与错误 key；\method{rawTcpFetchAndCommitTokensFenceInvalidOwnersWithoutChangingOffsets} 检查真实 TCP 的旧/future/未知/非 owner token；\method{groupConsumerPreservesRetainedPositionDropsRevokedPartitionsAndRecoversAfterExpiryAndRestart} 检查刷新、接管与重启恢复；\method{groupConsumerClosesItsOwnedRpcClientAfterPollingARecord} 检查资源归属。}
\lessonhints{请求校验必须同时核对 member、generation 和 partition owner；只比较 group 名称不足以围栏。}{重平衡后对每个分区分类：保留的本地 position、刚取得所有权需从 committed/0 开始的位置、已撤销必须丢弃的位置。}{沿用示例让 a 用旧 generation 提交较大 offset；检查错误码及所有相关 OffsetKey 是否保持原值。}
\lessonquestions{generation 围栏可以拒绝什么，为什么不能保证外部副作用 exactly-once？}{若 validate 解锁后再 commit，两个操作之间可能发生什么？}
\paragraph{自检解释。}服务端能拒绝旧 token 的 fetch/commit，却无法撤回已交付的记录或外部系统已执行的动作；若 validate 与追加不在同一 coordinator 临界区，rebalance 可插入两者之间，让旧 member 越权写进度。
\pitfalls{先校验、释放锁、再提交会留下 rebalance 竞态。把 \method{ILLEGAL_GENERATION} 和 \method{UNKNOWN_MEMBER} 都转成自动重新 join 会掩盖不同状态；真实 Kafka 组协议和事务提交更复杂，本课程单进程 coordinator 与本地 offset store 不等于高可用 \pathcode{__consumer_offsets}。}
\lessoncommand{20}
